\begin{abstract}
Swarms of language-model agents work together for months. Researchers must explain such swarms from their logs, and operators must control them. We ask which physics models describe one swarm, the AI Village, well enough for both. Over 51 goal periods and three scaffold regimes, we tested 16 physics models with 132 pre-registered hypotheses. Each hypothesis had to beat its strongest null and remove four impostors: the scheduler, external drives, shared model priors and simultaneous convergence. Four models are useful and fit one picture. Each agent is a kinetic spin that changes state only at its own calls. Agents couple only through messages that a call reads, mostly messages that name the reader. The coupling is weak: one message causes about 0.13 extra messages, and the talk fluctuations agree with no free parameter. Strong external fields set most of the state: the schedule, the goal text, the rooms and agent style. Each agent relaxes as a fast kick of about 7 calls on a slow well of about 100 calls. Few collective modes rise above the noise floor. We tested and mostly dropped the superagent and information-dynamics lines of inquiry. From the picture we derive 16 ranked observables for operators. Tests on reserved data have confirmed almost no result yet.
\end{abstract}
